\documentclass[11pt,a4paper]{article}
\usepackage{amsmath,amssymb}
\usepackage{listings}
\usepackage{hyperref}
\usepackage[margin=1in]{geometry}
\usepackage{fancyhdr}

\pagestyle{fancy}
\fancyhf{}
\fancyhead[L]{Module 16: Network Programming}
\fancyhead[R]{C Programming Course}
\fancyfoot[C]{\thepage}

\lstset{
  language=C,
  basicstyle=\ttfamily\small,
  keywordstyle=\bfseries,
  commentstyle=\itshape,
  numbers=left,
  numberstyle=\tiny,
  frame=single,
  breaklines=true,
  tabsize=4
}

\title{Module 16: Network Programming}
\author{C Programming Course}
\date{}

\begin{document}
\maketitle
\thispagestyle{fancy}

\section{Socket Programming Basics}
A socket is an endpoint for network communication. The typical server flow:
\texttt{socket} $\to$ \texttt{bind} $\to$ \texttt{listen} $\to$ \texttt{accept}.

\section{TCP/IP Communication}
\begin{lstlisting}
#include <stdio.h>
#include <string.h>
#include <unistd.h>
#include <arpa/inet.h>

int main(void) {
    int sockfd = socket(AF_INET, SOCK_STREAM, 0);
    if (sockfd < 0) { perror("socket"); return 1; }

    struct sockaddr_in addr;
    memset(&addr, 0, sizeof(addr));
    addr.sin_family = AF_INET;
    addr.sin_port = htons(8080);
    addr.sin_addr.s_addr = INADDR_ANY;

    if (bind(sockfd, (struct sockaddr *)&addr,
             sizeof(addr)) < 0) {
        perror("bind"); close(sockfd); return 1;
    }
    listen(sockfd, 5);
    printf("Server listening on port 8080\n");

    int client = accept(sockfd, NULL, NULL);
    if (client >= 0) {
        const char *msg = "Hello from server\n";
        write(client, msg, strlen(msg));
        close(client);
    }
    close(sockfd);
    return 0;
}
\end{lstlisting}

\section{UDP Programming}
UDP is connectionless. Use \texttt{SOCK\_DGRAM} instead of \texttt{SOCK\_STREAM},
and \texttt{sendto}/\texttt{recvfrom} instead of \texttt{read}/\texttt{write}.

\section{HTTP Client/Server}
An HTTP request is a text-based protocol over TCP. A minimal response:
\begin{lstlisting}
const char *response =
    "HTTP/1.1 200 OK\r\n"
    "Content-Type: text/plain\r\n"
    "Content-Length: 13\r\n"
    "\r\n"
    "Hello, World!";
write(client_fd, response, strlen(response));
\end{lstlisting}

\section{Multi-threaded Servers}
Use \texttt{pthread\_create} to handle each client in a separate thread.

\begin{lstlisting}
#include <pthread.h>

void *handle_client(void *arg) {
    int fd = *(int *)arg;
    /* read/write with client */
    close(fd);
    return NULL;
}
\end{lstlisting}

\section{Key Concepts}
\begin{itemize}
  \item Byte order: use \texttt{htons}/\texttt{ntohs} for port numbers.
  \item Always check return values of socket operations.
  \item Use \texttt{SO\_REUSEADDR} to avoid ``address already in use'' errors.
\end{itemize}

\end{document}
